\SourceLeaf{371}%
% A544-P0371-D01
Es ist also
\[
R+R'+R''+R'''=2(\alpha\beta abAB)-2(b)+2(B)+2(\beta)
\]
folglich
\[
\begin{aligned}
\Theta+\Theta'+\Theta''+\Theta'''={}&-g+G+\tfrac12(\alpha\beta abAB)\\
&+\tfrac12(\beta)-\tfrac12(b)+\tfrac12(B)\\
={}&-g+G+S
\end{aligned}
\]

% A544-P0371-D02
\begingroup\small\setlength{\arraycolsep}{2.6pt}
\[
\begin{array}{ccccc||ccc||ccc||ccc}
\multicolumn{2}{c}{ab}&\multicolumn{3}{c||}{+\quad+}&\multicolumn{3}{c||}{-\quad+}&\multicolumn{3}{c||}{-\quad-}&\multicolumn{3}{c}{+\quad-}\\\hline
A&B&\alpha&\beta&S&\alpha&\beta&S&\alpha&\beta&S&\alpha&\beta&S\\
+&+&+&+&+1&-&+&+1&-&-&+1&+&-&+1\\
&&+&-&-1&+&+&0&-&+&+1&-&-&0\\
-&+&+&-&0&+&+&+1&-&+&+2&-&-&+1\\
&&-&-&-1&+&-&-1&+&+&+1&-&+&+1\\
-&-&-&-&-1&+&-&-1&+&+&+1&-&+&+1\\
&&-&+&-1&-&-&-2&+&-&-1&+&+&0\\
+&-&-&+&0&-&-&-1&+&-&0&+&+&+1\\
&&+&+&-1&-&+&-1&-&-&-1&+&-&-1\\
\end{array}
\]
\endgroup

% A544-P0371-D03
\Article{14}
Hienach bekommt nun die erste Regel folgende Gestalt:
\[
\begin{aligned}
4\,\text{Dec.}={}&\text{I.}\;-4\Sigma\varepsilon\ \text{von denen, wo }y\text{ ganz, }[x]\text{ gerade}\\
&\text{II.}\;+4\Sigma\varepsilon\ \text{von denen, wo }[x]\text{ gerade, }[y]\text{ ungerade}\\
&\text{III.}\;-\Sigma\varepsilon\ \text{von allen}\\
&\text{IV.}\;+4\Sigma\varepsilon\ \text{von denen, wo nicht zugleich}\\
&\qquad [x]\text{ und }[y]\text{ gerade}\\
&+Q+S
\end{aligned}
\]

% A544-P0371-D04
In unserm Beispiel
\[
\begin{array}{rc}
\mathrm{I}\dotsfour&0\\
\mathrm{II}\dotsfour&-4\\
\mathrm{III}&+1\\
\mathrm{IV}&0\\
Q&-5\\
S&0\\\cline{2-2}
&-8
\end{array}
\]
Man denke sich nun in III diejenigen besonders bemerkt, wo $y$ ganz, $[x]$ gerade, so ist

\SourceLeaf{372}%
% A544-P0372-D01
\[
\begin{gathered}
\left.\begin{array}{rl}
\mathrm{III.}&\Sigma\varepsilon\ \text{von allen}\\
&-4\Sigma\varepsilon\ \text{der besonderen}
\end{array}\right\}\\
={}\text{Intensor }(\tfrac12-\omega)m-\text{Intens. }\omega m=-W
\end{gathered}
\]
Hier ist
\[
W=\begin{array}{cc|rr}
a&b&&\\
+&+&-3&-1\\
-&+&-2&0\\
-&-&+2&0\\
-&+&+3&+1\\[3pt]
&&\frac{m-1}{2}&\frac{m-1}{2}\\
&&\text{par}&\text{impar}
\end{array}
\]

% A544-P0372-D02
Also
\[
\begin{aligned}
4\,\text{Decident}={}&\text{I.}\;-4\Sigma\varepsilon\quad y\text{ ganz, }[x]\text{ gerade}\\
&\text{II.}\;+4\Sigma\varepsilon\quad[x]\text{ gerade, }[y]\text{ ungerade}\\
&\text{III.}\;-4\Sigma\varepsilon\quad y\text{ ganz, }[x]\text{ gerade}\\
&\text{IV.}\;+4\Sigma\varepsilon\quad\text{alle wo nicht zugleich }[x],[y]\text{ gerade}\\
&+Q+S+W
\end{aligned}
\]

% A544-P0372-D03
\begin{center}Tabelle für $\frac14(Q+S+W)$\end{center}
\begingroup\small\setlength{\arraycolsep}{2pt}
\[
\begin{array}{cccccc|rr||cccccc|rr}
a&b&A&B&\alpha&\beta&&&a&b&A&B&\alpha&\beta&&\\
+&+&+&+&+&+&0&0&-&-&+&+&-&-&0&0\\
&&&&+&-&-1&0&&&&&-&+&+1&0\\
&&-&+&+&-&-1&-1&&&-&+&-&+&0&-1\\
&&&&-&-&-1&-1&&&&&+&+&0&-1\\
&&-&-&-&-&-2&-1&&&-&-&+&+&0&-1\\
&&&&-&+&-1&-1&&&&&+&-&-1&-1\\
&&+&-&-&+&-1&0&&&+&-&+&-&0&0\\
&&&&+&+&-1&0&&&&&-&-&0&0\\[3pt]
-&+&+&+&-&+&0&0&+&-&+&+&+&-&+1&+1\\
&&&&+&+&0&0&&&&&-&-&0&0\\
&&-&+&+&+&-1&-1&&&-&+&-&-&+1&0\\
&&&&+&-&-1&-1&&&&&-&+&0&-1\\
&&-&-&+&-&-2&-1&&&-&-&-&+&+1&0\\
&&&&-&-&-2&-1&&&&&+&+&0&-1\\
&&+&-&-&-&-1&0&&&+&-&+&+&+1&+1\\
&&&&-&+&-1&0&&&&&+&-&0&0\\
\end{array}
\]
\endgroup

\SourceLeaf{373}%
% A544-P0373-D01
\Article{15}
Die zweite Methode ist folgende:
\[
\begin{aligned}
\text{Decident}={}&\text{I.}\;+\Sigma\varepsilon,\ \text{wo }Y\text{ ganz, }[X]\text{ gerade}\\
&-4\Sigma\varepsilon,\ \text{unter diesen, wo noch }y\text{ ganz, }[x]\text{ gerade}\\
&\text{II.}\;+\Sigma\lambda\varepsilon,\ \text{wo }Y\text{ ganz, }[X]\text{ gerade;}\\
&\qquad\lambda\text{ ist der Intensor von }p\\
&\text{II.}\;-\Sigma\lambda'\varepsilon,\ \text{wo }X\text{ ganz, }[Y]\text{ ungerade}\\
&\qquad\lambda'\text{ der Intensor von }ip=1,\,2,\,3,\,0\\
&\qquad\qquad\text{wenn }\lambda=0,\,1,\,2,\,3\\
&\text{IV.}\;+\Sigma\lambda\varepsilon,\ \text{wo }Y\text{ ganz, }[X]\text{ gerade,}\\
&\qquad\lambda\text{ der Intensor von }p\\
&\text{IV.}\;+\Sigma\lambda'\varepsilon,\ \text{wo }X\text{ ganz, }[Y]\text{ gerade}\\
&\qquad\lambda'\text{ der Intensor von }im-ip=0\quad3\quad2\quad1\\
&\qquad\qquad\text{wenn Int. }p=0\quad1\quad2\quad3\\
&+q
\end{aligned}
\]

% A544-P0373-D02
Hier ist $q=0$, wenn $\frac{M-1}{2}$ ungerade i.\,e. nur durch $1+i$, nicht durch $2$ theilbar, und nicht zugleich $AB\,+\,-$, hingegen übrigens
\begingroup\small\setlength{\arraycolsep}{2pt}
\[
\begin{array}{cc|l||rr|rr|rr|rr}
&&&\alpha&\beta&\alpha&\beta&\alpha&\beta&\alpha&\beta\\
A&B&&+&+&-&+&-&-&+&-\\
+&+&\text{Int. }\frac12(m-\omega\mu)&+3&+1&+3&+1&+0&+2&+0&+2\\
-&+&0&\multicolumn{2}{c|}{0}&\multicolumn{2}{c|}{0}&\multicolumn{2}{c|}{0}&\multicolumn{2}{c}{0}\\
-&-&-\text{Int. }\frac12(m+\omega\mu)&-0&-2&-0&-2&-3&-1&-3&-1\\
+&-&-\text{Int. }\omega\mu&\multicolumn{2}{c|}{-0}&\multicolumn{2}{c|}{-1}&\multicolumn{2}{c|}{-2}&\multicolumn{2}{c}{-3}
\end{array}
\]
\endgroup
wo doppelte Zahlen stehen, gilt die erste für gerade $\frac{m-1}{2}$, die andere für ungerade.

% A544-P0373-D03
In unserm Beispiele:
\[
\begin{array}{rc|r|r|r|r}
\mathrm{I.}&y&Y&20x&20Y&\varepsilon\\
&0&-1&+13&+9&+1^{*}\\
&0&-2&+26&+18&+1\\
&+1&-1&+22&+46&+1\\
&&&&&+3(+1)
\end{array}
\qquad
\begin{array}{r}
+3\\-4\\-1\\\hline
\text{Dec.}=-2
\end{array}
\]

% A544-P0373-D04
II. desunt.
\begingroup\small\setlength{\arraycolsep}{2pt}
\[
\begin{array}{rc|r|r|r|r|r|r|r}
\mathrm{IV.}&X&12x&12y&12Y&\varepsilon&\lambda&\lambda'&\lambda'\varepsilon\\
&0&+20+20\omega&-9&-37&-1&2&2&-2\\
&+1&+24+20\omega&-3&-39&+1&3&1&+1\\
&+2&+28+20\omega&+3&-41&-1&0&0&0\\
&&&&&&&&-1
\end{array}
\]
\endgroup

\SourceLeaf{374}%
% A544-P0374-D01
Was aus II genommen ist, vereinigt sich in folgendes Resultat
\[
\begin{aligned}
\mathrm{II.}\;&-\Sigma\varepsilon,\ \text{wo }Y\text{ ganz, }[X]\text{ gerade}\\
&+4\Sigma\varepsilon,\ \text{von eben diesen, wenn zugleich }[x]\text{ gerade,}\\
&\qquad[y]\text{ ungerade}\\
\mathrm{II.}\;&+\Sigma\lambda'\varepsilon,\ \text{wo }Y\text{ ganz, }[X]\text{ gerade}\\
&-\Sigma\lambda'\varepsilon,\ \text{wo }X\text{ ganz, }[Y]\text{ ungerade,}\\
&\qquad\lambda'\text{ der Intensor von }ip
\end{aligned}
\]

% A544-P0374-D02
Die beiden letzten Theile vereinigen sich wiederum in
\[
\begin{aligned}
\mathrm{III.}\;&+\Sigma\varepsilon,\ \text{wo }[X]\text{ gerade, }[Y]\text{ ungerade}\\
&-4\Sigma\varepsilon,\ \text{wenn zugleich }x\text{ eine Ganze, }[y]\text{ ungerade}\\
&+r+s
\end{aligned}
\]
wo $r=0$, wenn $AB\dotsfive\left\{\begin{array}{cc}+&+\\-&+\\-&-\end{array}\right.$

% A544-P0374-D03
und $r=\text{Int. }\omega mi$, wenn $AB\dotstwo+\,-$
\[
\begin{gathered}
s=-\text{Int. }(\tfrac12-\omega)mi,\quad\text{wenn }\frac{M-1}{2}\text{ gerade und }B\text{ positiv,}\\
\text{in allen andern Fällen }=0
\end{gathered}
\]

% A544-P0374-D04
In unserm Beispiele
\[
\text{III. Fällt aus.}\quad r=0,\ s=0
\]
Was aus IV genommen ist, vereinigt sich in folgendes Resultate
\[
\begin{aligned}
\mathrm{IV.}\;&+4\Sigma\varepsilon,\ \text{wo }Y\text{ ganz, }[X]\text{ gerade und nicht zugleich}\\
&\qquad[x]\text{ gerade, }[y]\text{ gerade}\\
&-\Sigma\lambda'\varepsilon,\ \text{wo }Y\text{ ganz, }[X]\text{ gerade}\\
&+\Sigma\lambda'\varepsilon,\ \text{wo }X\text{ ganz, }[Y]\text{ gerade}\\
&\qquad\text{wo }\lambda'\text{ Int. von }im-ip
\end{aligned}
\]

% A544-P0374-D05
Die zwei letzten Theile vereinigen sich wiederum zu
\[
\begin{aligned}
\mathrm{V.}\;&+\Sigma\varepsilon,\ \text{wo nicht zugleich }[X]\text{ und }[Y]\text{ gerade}\\
&-4\Sigma\varepsilon,\ \text{wo zugleich }x\text{ eine Ganze, }[y]\text{ gerade}\\
&+w
\end{aligned}
\]
Hier ist $w=0$, wenn $\frac{M-1}{2}$ gerade und $A$ positiv; in allen übrigen Fällen
\[
=-\text{Intensor }(\tfrac12i+\omega)m
\]

\SourceLeaf{375}%
% A544-P0375-D01
In unserm Beispiele
\[
\begin{gathered}
\mathrm{V.}\\
\begin{array}{c|r|r|r|r}
y&6x&6X&6Y&\varepsilon\\
0&+13&+9&-20&+1
\end{array}\\
w=-2
\end{gathered}
\]

% A544-P0375-D02
\begin{center}Tafel für $q,r,s,w$ und deren Summe.\end{center}
\begingroup\scriptsize\setlength{\arraycolsep}{1.1pt}\renewcommand{\arraystretch}{1.12}
\[
\begin{array}{cccccc|rrrrr|rrrrr|rrrrr|rrrrr}
a&b&A&B&\alpha&\beta&\multicolumn{5}{c|}{\begin{array}{c}\frac{M-1}{2}\ \text{ger.}\\\frac{m-1}{2}\ \text{ger.}\end{array}}&\multicolumn{5}{c|}{\begin{array}{c}\frac{M-1}{2}\ \text{ung.}\\\frac{m-1}{2}\ \text{ger.}\end{array}}&\multicolumn{5}{c|}{\begin{array}{c}\frac{M-1}{2}\ \text{ger.}\\\frac{m-1}{2}\ \text{ung.}\end{array}}&\multicolumn{5}{c|}{\begin{array}{c}\frac{M-1}{2}\ \text{ung.}\\\frac{m-1}{2}\ \text{ung.}\end{array}}\\[4pt]
+&+&+&+&+&+&+3&0&0&0&+3&0&0&0&0&0&+1&0&-2&0&-1&0&0&0&-2&-2\\
&&&&+&-&0&0&0&0&0&0&0&0&0&0&+2&0&-2&0&0&0&0&0&-2&-2\\
&&-&+&+&-&0&0&0&0&0&0&0&0&0&0&0&0&-2&-2&-4&0&0&0&-2&-2\\
&&&&-&-&0&0&0&0&0&0&0&0&0&0&0&0&-2&-2&-4&0&0&0&-2&-2\\
&&-&-&-&-&-3&0&0&0&-3&0&0&0&0&0&-1&0&0&-2&-3&0&0&0&-2&-2\\
&&&&-&+&0&0&0&0&0&0&0&0&0&0&-2&0&0&-2&-4&0&0&0&-2&-2\\
&&+&-&-&+&-1&+1&0&0&0&-1&+1&0&0&0&-1&+1&0&0&0&-1&+1&0&-2&-2\\
&&&&+&+&0&+1&0&0&+1&0&+1&0&0&+1&0&+1&0&0&+1&0&+1&0&-2&-1\\[4pt]
-&+&+&+&-&+&+3&0&0&0&+3&0&0&0&-1&-1&+1&0&-2&0&-1&0&0&0&-3&-3\\
&&&&+&+&+3&0&0&0&+3&0&0&0&-1&-1&+1&0&-2&0&-1&0&0&0&-3&-3\\
&&-&+&+&+&0&0&0&-1&-1&0&0&0&-1&-1&0&0&-2&-3&-5&0&0&0&-3&-3\\
&&&&+&-&0&0&0&-1&-1&0&0&0&-1&-1&0&0&-2&-3&-5&0&0&0&-3&-3\\
&&-&-&+&-&-3&0&0&-1&-4&0&0&0&-1&-1&-1&0&0&-3&-4&0&0&0&-3&-3\\
&&&&-&-&-3&0&0&-1&-4&0&0&0&-1&-1&-1&0&0&-3&-4&0&0&0&-3&-3\\
&&+&-&-&-&-2&+2&0&0&0&-2&+2&0&-1&-1&-2&+2&0&0&0&-2&+2&0&-3&-3\\
&&&&-&+&-1&+2&0&0&+1&-1&+2&0&-1&0&-1&+2&0&0&+1&-1&+2&0&-3&-2\\[4pt]
-&-&+&+&-&-&0&0&-1&0&-1&0&0&0&-1&-1&+2&0&-3&0&-1&0&0&0&-3&-3\\
&&&&-&+&+3&0&-1&0&+2&0&0&0&-1&-1&+1&0&-3&0&-2&0&0&0&-3&-3\\
&&-&+&-&+&0&0&-1&-1&-2&0&0&0&-1&-1&0&0&-3&-3&-6&0&0&0&-3&-3\\
&&&&+&+&0&0&-1&-1&-2&0&0&0&-1&-1&0&0&-3&-3&-6&0&0&0&-3&-3\\
&&-&-&+&+&0&0&0&-1&-1&0&0&0&-1&-1&-2&0&0&-3&-5&0&0&0&-3&-3\\
&&&&+&-&-3&0&0&-1&-4&0&0&0&-1&-1&-1&0&0&-3&-4&0&0&0&-3&-3\\
&&+&-&+&-&-3&+3&0&0&0&-3&+3&0&-1&-1&-3&+3&0&0&0&-3&+3&0&-3&-3\\
&&&&-&-&-2&+3&0&0&+1&-2&+3&0&-1&0&-2&+3&0&0&+1&-2&+3&0&-3&-2\\[4pt]
+&-&+&+&+&-&0&0&-1&0&-1&0&0&0&0&0&+2&0&-3&0&-1&0&0&0&-2&-2\\
&&&&-&-&0&0&-1&0&-1&0&0&0&0&0&+2&0&-3&0&-1&0&0&0&-2&-2\\
&&-&+&-&-&0&0&-1&0&-1&0&0&0&0&0&0&0&-3&-2&-5&0&0&0&-2&-2\\
&&&&-&+&0&0&-1&0&-1&0&0&0&0&0&0&0&-3&-2&-5&0&0&0&-2&-2\\
&&-&-&-&+&0&0&0&0&0&0&0&0&0&0&-2&0&0&-2&-4&0&0&0&-2&-2\\
&&&&+&+&0&0&0&0&0&0&0&0&0&0&-2&0&0&-2&-4&0&0&0&-2&-2\\
&&+&-&+&+&0&0&0&0&0&0&0&0&0&0&0&0&0&0&0&0&0&0&-2&-2\\
&&&&+&-&-3&0&0&0&-3&-3&0&0&0&-3&-3&0&0&0&-3&-3&0&0&-2&-5\\
\end{array}
\]
\endgroup

